\chapter*{Glossary}
\addcontentsline{toc}{chapter}{Glossary}

\begin{description}[Context window]
\item[Attention] A content-dependent mixing operation that forms query, key, and value representations and aggregates values according to query-key compatibility.
\item[Context window] The maximum number of tokens a model can condition on during a forward pass.
\item[Instruction tuning] Fine-tuning a pretrained model on examples that pair user-like instructions with target responses.
\item[KV cache] Stored key and value tensors from previous decoding steps, reused to avoid recomputing attention over the prefix.
\item[Preference model] A model or implicit objective that represents which of two candidate responses should be preferred under a target policy or annotation protocol.
\item[Retrieval-augmented generation] A generation pattern in which external documents are retrieved and inserted into the model context before or during answer generation.
\item[Test-time compute] Additional computation spent at inference time through longer reasoning traces, multiple samples, verification, search, or refinement.
\end{description}
